\chapter{Onychomycosis (Tinea Unguium)}
\index{Onychomycosis (Tinea Unguium)}
\label{sec:Onychomycosis (Tinea Unguium)}

\section{Overview}
Onychomycosis is a fungal infection of the nail unit, most commonly caused by dermatophytes, particularly Trichophyton rubrum. It accounts for approximately 50\% of all nail disorders and affects 10--20\% of the general population. The condition is more common in older adults, diabetics, immunocompromised individuals, and those with peripheral vascular disease.
\section{Classification}

\begin{description}[font=\normalfont\bfseries,leftmargin=3em,labelwidth=2.5em]
  \item[Cause] Infectious (Fungal)
  \item[Organ System] Skin
  \item[Pathophysiology] Dermatophyte invasion of the nail plate via the hyponychium or nail bed
  \item[Duration] Chronic (months to years)
  \item[Onset] Gradual
  \item[Mode of Transmission] Direct contact / environmental
  \item[Clinical Course] Chronic, progressive without treatment
  \item[Severity] Mild to Moderate
  \item[Extent of Disease] Localized
  \item[Age Group] Adults (increasing with age)
  \item[Epidemiology] Prevalence 10--20\%; higher in diabetics, elderly, athletes
  \item[ICD-11 Code] 1F28.0
\end{description}

\section{Causes}
Dermatophytes (Trichophyton rubrum, Trichophyton mentagrophytes) cause 90\% of cases. Yeasts (Candida albicans) and non-dermatophyte molds may also be responsible. The fungus invades the nail plate through the distal or lateral nail bed, spreading proximally and causing progressive nail destruction. Warm, moist environments and shared public facilities promote transmission.

\section{Risk Factors}

\begin{itemize}[noitemsep]
  \item Advancing age
  \item Diabetes mellitus
  \item Peripheral vascular disease
  \item Immunosuppression
  \item Occlusive footwear and excessive sweating
  \item Use of public swimming pools, gyms, and showers
  \item Previous nail trauma
  \item Family history of onychomycosis
  \item Tinea pedis (athlete's foot)
\end{itemize}

\section{Signs and Symptoms}

\subsection{Early symptoms}
\begin{itemize}[noitemsep]
  \item Early manifestations of onychomycosis (tinea unguium) may be subtle and nonspecific
\end{itemize}

\subsection{Common symptoms}
\begin{itemize}[noitemsep]
  \item \subsection{Common symptoms}
  \item \textbf{Common symptoms:}
  \item Thickened, discolored (white, yellow, brown) nail plate
  \item Subungual debris and hyperkeratosis
  \item Nail plate separation from the nail bed (onycholysis)
  \item Nail plate crumbling and destruction
  \item Foul odor from accumulated debris
  \item Pain with pressure from thickened nail
  \item Difficulty performing daily activities
\end{itemize}

\subsection{Less common symptoms}
\begin{itemize}[noitemsep]
  \item Atypical or less common presentations of onychomycosis (tinea unguium) may occur
\end{itemize}

\subsection{Emergency warning signs}
\begin{itemize}[noitemsep]
  \item Severe or worsening symptoms of onychomycosis (tinea unguium) require immediate medical evaluation
\end{itemize}
\section{Progression and Stages}
Untreated onychomycosis progresses slowly over years. The infection typically begins at the distal or lateral edge of the nail and spreads proximally. Total nail plate involvement can develop over 1--5 years. Secondary bacterial infection may occur, and the nail may become completely destroyed. Spontaneous resolution is rare.

\section{Complications}

\begin{itemize}[noitemsep]
  \item Secondary bacterial cellulitis, particularly in diabetics
  \item Diabetic foot ulcers and complications
  \item Permanent nail dystrophy and disfigurement
  \item Spread to other nails or skin
  \item Chronic pain and difficulty walking
  \item Reduced quality of life
  \item Emotional and psychological effects
\end{itemize}

\section{Diagnosis}

\begin{itemize}[noitemsep]
  \item Clinical examination with characteristic nail findings
  \item Potassium hydroxide (KOH) preparation with fungal culture
  \item Periodic acid-Schiff (PAS) staining of nail clippings (most sensitive)
  \item Dermatophyte test medium (DTM) culture
  \item PCR for rapid species identification
  \item Fungal culture to identify causative organism
\end{itemize}

\section{Prevention}

\begin{itemize}[noitemsep]
  \item Maintain good foot hygiene and keep feet dry
  \item Wear breathable footwear and moisture-wicking socks
  \item Avoid walking barefoot in public facilities
  \item Disinfect nail care instruments
  \item Treat associated tinea pedis
  \item Go for regular check-ups so any problems are caught early
\end{itemize}

\section{Treatment Options}

\begin{itemize}[noitemsep]
  \item Topical antifungal lacquers (ciclopirox, amorolfine) for mild distal disease
  \item Oral antifungal therapy (terbinafine, itraconazole) for moderate to severe disease
  \item Terbinafine 250 mg daily for 6--12 weeks (or toenail: 12 weeks)
  \item Laser therapy and photodynamic therapy (emerging options)
  \item Nail debridement and chemical avulsion
  \item Surgical nail removal for refractory cases
  \item Lifestyle changes and self-care
\end{itemize}

\section{Living with It}

\begin{itemize}[noitemsep]
  \item Follow your treatment plan as prescribed
  \item Keep up with regular appointments with your healthcare provider
  \item Adjust your daily habits as needed
  \item Reach out to family, friends, or support groups for help
  \item Keep learning about your condition and how to manage it
\end{itemize}

\section{Outlook}
With appropriate oral antifungal therapy, cure rates of 70--80\% for toenail infections are achievable. However, recurrence rates are high (10--50\%) even after successful treatment. Fingernail infections respond better than toenail infections. New nail growth takes 6--12 months for fingernails and 12--18 months for toenails to appear normal. Long-term prevention strategies are important to minimize recurrence.
